\section{논의}

\subsection{결과 요약}
이 연구는 화면 테두리를 잠깐 밝히는 시각 경고 신호가, 개인 기기에서 하는 웹 과제에서도 같은 사람의 단순 반응시간을 줄이는지 확인했다.
사전에 정한 규칙에 따라 가설은 \ResVerdict.

% 아래 두 갈래의 해석 틀은 데이터 수집 전에 써 두었다(Git 기록 참고). 결과가 나온 뒤 고르는 것이 아니라 미리 정한 대로 따른다.
\ifnum\ResSupported=1
\subsection{해석}
결과는 실험실 연구의 각성 효과 \citep{weinbach2012,fan2002}와 같은 방향이다.
이 연구의 기여는 효과의 존재가 아니라 환경에 있다.
참가자는 서로 다른 기기와 화면 크기, 입력 방식으로, 감독 없이 실험에 참여했다.
그런 조건에서도 UI 장식에 가까운 테두리 점등이 반응을 빠르게 했다.
따라서 ``곧 나온다''는 시각 신호는 실제 사용 환경에서도 반응 준비를 돕는다고 볼 수 있다.
\else
\subsection{해석}
이 결과를 ``경고 신호는 효과가 없다''로 읽으면 안 된다.
가능한 설명은 미리 네 가지로 정해 두었고, 원자료로 확인할 수 있는 것부터 확인한다.
(1) \emph{표본 부족}: 이 표본이 잡아낼 수 있는 최소 효과크기($d_z\approx\MDEdz$)보다 실제 효과가 작을 수 있다.
(2) \emph{신호를 못 봄}: 테두리는 화면 가장자리에 있어, 휴대폰을 쥔 손이나 시선 위치에 따라 눈에 덜 띌 수 있다. 입력 방식별 차이로 일부 확인한다.
(3) \emph{시간 정보 효과의 상쇄}: 신호 없음 조건에서도 기다릴수록 목표가 곧 나온다는 예측이 생겨 반응이 빨라진다 \citep{niemi1981}. 이 효과가 두 조건의 차이를 줄였을 수 있다.
(4) \emph{성급한 반응과의 교환}: 신호가 반응을 앞당기는 대신 조기 반응을 늘렸다면, 그 시행은 분석에서 빠져 효과가 작게 보인다. 조건별 조기 반응 건수로 확인한다.
\fi

\TODO{데이터를 본 뒤 작성: 위 설명 가운데 원자료로 확인한 것과 그 결과(T10-C36). 가설을 고쳤다면 고치기 전 문장과 고친 문장, 고친 이유(T10-C37). 고치지 않았다면 그 이유.}

\subsection{한계}\label{sec:limits}
\begin{enumerate}
\item \textbf{메커니즘을 분리하지 못한다.} 경고 신호는 각성을 높이는 동시에 목표가 정확히 500\,ms 뒤에 나온다는 시간 정보도 준다.
이 두 과정은 대부분의 설계에서 섞여 있고 \citep{weinbach2012}, 이 실험도 둘을 구분하지 못한다.
따라서 이 연구가 말할 수 있는 것은 ``이 신호가 반응시간을 바꿨는가''까지이며, 그 이유가 각성인지 시간 예측인지는 말할 수 없다.
\item \textbf{하드웨어 지연은 측정하지 않았다.} 기기마다 반응시간이 수십 ms 이상 과대 측정될 수 있다 \citep{pronk2020}.
참가자 내 차이는 이 고정 지연을 상쇄하지만, 이 논문의 \emph{절대} 반응시간을 다른 연구와 직접 비교해서는 안 된다.
\item \textbf{표본.} 지인과 커뮤니티 이용자의 편의 표본이고, 인원은 검정력 분석이 아니라 수집 가능 범위로 정했다. 결과를 일반 인구로 넓히기 어렵다.
\item \textbf{감독 없는 환경.} 참가자의 주의 상태, 주변 소음, 화면 밝기를 통제하지 못했다. 탭 전환은 감지해 제외했지만 화면 밖의 방해는 감지할 수 없다.
\item \textbf{재참여 판별.} 재참여는 같은 브라우저의 저장소로만 판별하므로, 다른 기기나 시크릿 창으로 다시 참여한 경우는 걸러내지 못한다.
\item \textbf{자극 크기.} 테두리 두께와 원 크기는 CSS 픽셀로 고정했다. 화면 크기에 따라 실제 시각 크기와 눈에 띄는 정도가 달랐다.
\end{enumerate}

\subsection{결론}
\TODO{데이터를 본 뒤 작성: 한 문단 이상(T10-C38). 원자료와 위 분석에 있는 내용만 쓴다(T10-C40).}
